% !TeX root = ../../Book3.tex
% 纲要标题：### 21.5 Gauss–Green 公式与等周问题
% 纲要位置：几何学计划纲要.md，第 2998 行。
% 正文状态：初稿；尚未正式审核。

% 纲要内容开始
% 
% * 积分分部公式
% * BV 函数的切片与余面积公式
% * Euclid 等周不等式
% * 固定体积的周长极小化
% 
% ---
% 
% 纲要内容结束

\section{Gauss–Green 公式与等周问题}
\label{g3:sec:2105}

边界结构定理把周长测度落实为约化边界上的面积。现在可以沿两条方向使用它：把光滑区域的积分分部推广到有限周长集合；把周长作为能量，在可测集合之间比较大小。切片与余面积公式连接这两条方向，前者把一个方向的变化分解到直线上，后者把一般 BV 函数的变化分解到水平集上。

\subsection{积分分部公式}
\label{g3:subsec:2105:01}
本节仍取开集 \(\Omega\subset\mathbb R^n\)，\(n\ge1\)。外法向约定为
\[
	D\chi_E=-\nu_E\,|D\chi_E|.
\]
这个负号可用区间检查：\(D\chi_{(a,b)}=\delta_a-\delta_b\)，而左端点的外法向为 \(-1\)，右端点的外法向为 \(1\)。

\termidx[Gauss Green gongshi]{Gauss--Green 公式}
\foreignidx{Gauss--Green formula}
\begin{theorem}[有限周长集合的 Gauss--Green 公式]\label{g3:thm:finite-perimeter-gauss-green}
	设 \(\Omega\subset\mathbb R^n\) 开，\(n\ge1\)，\(E\subset\mathbb R^n\) 在 \(\Omega\) 中局部有限周长，\(\nu_E\) 为其约化边界上的广义外单位法向。对每个 \(X\in C_c^1(\Omega;\mathbb R^n)\)，有
	\begin{equation}\label{g3:eq:finite-perimeter-gauss-green}
		\int_E\operatorname{div}X\,\mathrm dx
		=\int_{\partial^*E\cap\Omega}X\cdot\nu_E\,\mathrm d\mathcal H^{n-1}.
	\end{equation}
	这里左端在 \(E\cap\Omega\) 上积分。更一般地，若 \(\varphi\in C^1(\Omega)\)，则
	\[
		\int_E\bigl(\varphi\operatorname{div}X+\nabla\varphi\cdot X\bigr)\,\mathrm dx
		=\int_{\partial^*E\cap\Omega}\varphi X\cdot\nu_E\,\mathrm d\mathcal H^{n-1}.
	\]
\end{theorem}
\begin{proof}
	分布导数的定义给出
	\[
		\int_E\operatorname{div}X\,\mathrm dx=-\int_\Omega X\cdot\mathrm dD\chi_E.
	\]
	这一定义先对光滑测试场成立，再将 \(C_c^1\) 测试场及其一阶导数一致逼近即可得到上述形式。由\cref{g3:thm:de-giorgi-structure}及\cref{g3:prop:one-dimensional-perimeter-structure}，
	\[
		D\chi_E=-\nu_E\,\mathcal H^{n-1}\!\llcorner(\partial^*E\cap\Omega).
	\]
	代入便得第一式；把测试场换成 \(\varphi X\)，并展开散度，得到第二式。紧支条件使所有积分只涉及一个有限周长的紧邻域。
\end{proof}

若 \(E\) 有界，\(\overline E\Subset\Omega\)，且 \(X\) 在 \(\overline E\) 的邻域内为 \(C^1\)，可取等于一的紧支截断函数，将公式用于截断后的场。这样便能使用 \(X(x)=x\) 等并不紧支的局部测试场。对于一般 \(u\in BV_{\mathrm{loc}}(\Omega)\)，相应公式仍是
\[
	\int_\Omega u\operatorname{div}X\,\mathrm dx=-\int_\Omega X\cdot\mathrm dDu.
\]
集合情形的额外内容在于：右端的向量测度已经有了外法向与边界面积的具体表示。

\begin{example}\label{g3:exm:annulus-gauss-green}
	设 \(n\ge2\)，\(0<r<R\)，\(E=B_R(0)\setminus\overline{B_r(0)}\)。外球面上的外法向是 \(x/R\)，内球面上的外法向是 \(-x/r\)。以等于 \(x\) 的截断场测试，得到
	\[
		n|E|=R\mathcal H^{n-1}(\partial B_R)-r\mathcal H^{n-1}(\partial B_r)
		=n\omega_n(R^n-r^n).
	\]
	内球面贡献负通量，是因为它的外侧指向被挖去的球。两部分对周长的贡献则都取正值。
\end{example}

\subsection{BV 函数的切片与余面积公式}
\label{g3:subsec:2105:02}
固定单位向量 \(e\)，写 \(x=y+te\)，其中 \(y\in e^\perp\)、\(t\in\mathbb R\)。令
\[
	\Omega_y=\{t:y+te\in\Omega\},\qquad
	u_y(t)=u(y+te),\qquad D_eu=e\cdot Du.
\]
这里 \(D_eu\) 是有符号测度，\(|D_eu|\) 是这个测度的全变差。一般有 \(|D_eu|\le|Du|\)，二者并不相同。

\begin{theorem}\label{g3:thm:bv-directional-slicing}
	设 \(\Omega\subset\mathbb R^n\) 开，\(n\ge1\)，\(u\in BV_{\mathrm{loc}}(\Omega)\)，\(e\in\mathbb R^n\) 为单位向量。对 \(y\in e^\perp\)，记 \(\Omega_y=\{t:y+te\in\Omega\}\)、\(u_y(t)=u(y+te)\)，并记方向导数测度 \(D_eu=e\cdot Du\)。对 \(\mathcal L^{n-1}\) 几乎所有 \(y\in e^\perp\)，有 \(u_y\in BV_{\mathrm{loc}}(\Omega_y)\)。若 \(B\subset\Omega\) 为 Borel 集，\(B_y=\{t:y+te\in B\}\)，则
	\begin{equation}\label{g3:eq:bv-directional-disintegration}
		|D_eu|(B)=\int_{e^\perp}|Du_y|(B_y)\,\mathrm d\mathcal L^{n-1}(y).
	\end{equation}
	对 \(\psi\in C_c(\Omega)\)，还有带符号的公式
	\[
		\int_\Omega\psi\,\mathrm dD_eu
		=\int_{e^\perp}\left(\int_{\Omega_y}\psi(y+te)\,\mathrm dDu_y(t)\right)
		\mathrm d\mathcal L^{n-1}(y).
	\]
	当 \(n=1\) 时，\(e^\perp\) 视为一个带单位质量的点，上述公式就是一维恒等式。
\end{theorem}
\begin{proof}
	旋转坐标后只需取 \(e=e_n\)。先固定长方体 \(Q=A\times I\Subset\Omega\)，其中 \(A\subset\mathbb R^{n-1}\) 为开长方体，\(I\) 为有界开区间。由 Fubini 定理，几乎所有 \(u_y\) 属于 \(L^1(I)\)。在
	\(\{\phi\in C_c^1(I):\|\phi\|_\infty\le1\}\) 中选取一个可数稠密族 \(\{\phi_j\}\)，包含零函数及每个函数的相反数；这里先按紧支集穷尽，再在 \(C^1\) 范数下取稠密族。令
	\[
		a_j(y)=\int_Iu_y(t)\phi_j'(t)\,\mathrm dt,
		\qquad V(y)=\sup_j a_j(y).
	\]
	每个 \(a_j\) 可测，所以 \(V\) 可测；由稠密性，\(V(y)\) 正是一维变差的对偶上确界，允许为无穷。

	证明 \(\int_AV\le|D_nu|(Q)\) 时，须让测试函数随 \(y\) 选择，不能直接交换上确界和积分。固定有限个 \(a_1,\ldots,a_N\)，按取得最大值的最小下标将 \(A\) 分为可测集 \(A_1,\ldots,A_N\)。令
	\(w(y)=\int_I|u_y(t)|\,\mathrm dt\)。用有限测度 \(w(y)\,\mathrm dy\) 的内外正则性，在各 \(A_j\) 内取两两不交的紧集，再用支集两两不交的光滑函数 \(0\le\eta_j\le1\) 逼近 \(\chi_{A_j}\)，使
	\(\sum_j\int_Aw|\eta_j-\chi_{A_j}|\)
	任意小。具体说，有限个紧集之间有正距离，可分别选不交的开邻域；先使遗漏的加权测度小，再缩小这些邻域即可。于是
	\[
		\Psi(y,t)=\sum_{j=1}^N\eta_j(y)\phi_j(t)\in C_c^1(Q),\qquad |\Psi|\le1.
	\]
	将它代入 \(D_nu\) 的对偶定义，并令加权逼近误差趋零，得到
	\[
		\int_A\max_{1\le j\le N}a_j(y)\,\mathrm dy\le|D_nu|(Q).
	\]
	令 \(N\to\infty\)，便有 \(\int_AV\le|D_nu|(Q)<\infty\)。因此几乎所有切片属于 \(BV(I)\)，且 \(V(y)=|Du_y|(I)\)。

	现在对 \(\Psi\in C_c^1(Q)\)，逐条直线应用一维分布导数定义，再用 Fubini 定理，得到
	\[
		-\int_Q u\partial_n\Psi\,\mathrm dx
		=\int_A\left(\int_I\Psi(y,t)\,\mathrm dDu_y(t)\right)\mathrm dy.
	\]
	右端绝对值不超过 \(\|\Psi\|_\infty\int_A|Du_y|(I)\,\mathrm dy\)。对 \(|\Psi|\le1\) 取上确界，得到反向不等式，故
	\[
		|D_nu|(A\times I)=\int_A|Du_y|(I)\,\mathrm dy.
	\]

	还需说明这不只是长方体的等式。对每个有理端点子区间 \(J\Subset I\)，同样的可数测试族说明 \(y\mapsto|Du_y|(J)\) 可测；开区间的递增穷尽和测度的单调类论证，把可测性推广到所有 Borel 子集。于是 \(|Du_y|\) 构成可测的测度族，
	\[
		\Lambda(B)=\int_A|Du_y|(B_y)\,\mathrm dy
	\]
	在 \(Q\) 上定义有限 Borel 测度。以上论证对所有较小长方体成立，故 \(\Lambda\) 与 \(|D_nu|\) 在生成 Borel 集的长方体族上一致，从而在所有 Borel 集上一致。带符号的公式先对光滑 \(\Psi\) 成立，再由这个变差控制及一致逼近推广到连续紧支函数。最后以可数个相对紧的有理长方体覆盖 \(\Omega\)；切片导数在重叠处由分布定义唯一，故可拼接，并在舍去一个共同的零测参数集后得到全部结论。
\end{proof}

\begin{corollary}\label{g3:cor:perimeter-line-slices}
	设 \(\Omega\subset\mathbb R^n\) 开，\(n\ge1\)，\(e\) 为单位向量，\(E\subset\mathbb R^n\) 在 \(\Omega\) 中局部有限周长。对 \(y\in e^\perp\) 记 \(\Omega_y=\{t:y+te\in\Omega\}\)、\(E_y=\{t:y+te\in E\}\)、\(B_y=\{t:y+te\in B\}\)，\(\nu_E\) 为广义外单位法向。则几乎所有直线切片 \(E_y\) 在 \(\Omega_y\) 内局部有限周长。对 Borel 集 \(B\subset\Omega\)，有
	\[
		\int_{e^\perp}\mathcal H^0(\partial^*E_y\cap B_y)\,\mathrm dy
		=\int_{\partial^*E\cap B}|\nu_E\cdot e|\,\mathrm d\mathcal H^{n-1}.
	\]
\end{corollary}
\begin{proof}
	在切片定理中取 \(u=\chi_E\)。由\cref{g3:prop:one-dimensional-perimeter-structure}，一维周长就是区间端点的计数测度；另一方面，
	\[
		D_e\chi_E=-(\nu_E\cdot e)\mathcal H^{n-1}\!\llcorner\partial^*E,
	\]
	所以其全变差的密度为 \(|\nu_E\cdot e|\)。代入切片公式即可。
\end{proof}

下一条公式按函数值切片。它既覆盖通常的梯度积分，也能记录跳跃面上的变化。

\begin{theorem}[BV 余面积公式]\label{g3:thm:bv-coarea}
	设 \(\Omega\subset\mathbb R^n\) 开，\(n\ge1\)，\(u\in BV_{\mathrm{loc}}(\Omega)\)，\(E_t=\{x\in\Omega:u(x)>t\}\)。则对几乎所有 \(t\in\mathbb R\)，\(E_t\) 在 \(\Omega\) 中局部有限周长；对每个 Borel 集 \(B\subset\Omega\)，有
	\begin{equation}\label{g3:eq:bv-coarea-measure}
		|Du|(B)=\int_{\mathbb R}|D\chi_{E_t}|(B)\,\mathrm dt
		=\int_{\mathbb R}\mathcal H^{n-1}(B\cap\partial^*E_t)\,\mathrm dt.
	\end{equation}
	积分允许为无穷。更一般地，对 \(u\in L^1_{\mathrm{loc}}(\Omega)\) 和开集 \(U\subset\Omega\)，把变差与周长按对偶上确界定义为扩展实数，则
	\begin{equation}\label{g3:eq:bv-coarea-open}
		V(u;U)=\int_{\mathbb R}P(E_t;U)\,\mathrm dt.
	\end{equation}
	因此，当 \(u\in L^1(\Omega)\) 时，右端有限也是 \(u\in BV(\Omega)\) 的充要条件。
\end{theorem}
\begin{proof}
	先在有界开集 \(U\Subset\Omega\) 上证明。周长的对偶定义可以只对一个可数稠密的测试场族取上确界，而每个
	\(t\mapsto\int_{E_t}\operatorname{div}X\)
	都由 Fubini 定理可测，故 \(t\mapsto P(E_t;U)\) 可测。

	第一步给出一个方向的不等式，不需要预先知道 \(u\) 属于 BV。对几乎所有 \(x\)，有层蛋糕表示
	\[
		u(x)=\int_{\mathbb R}\bigl(\chi_{E_t}(x)-\chi_{\{0>t\}}\bigr)\,\mathrm dt,
		\qquad
		\int_{\mathbb R}\bigl|\chi_{E_t}(x)-\chi_{\{0>t\}}\bigr|\,\mathrm dt=|u(x)|.
	\]
	若 \(X\in C_c^1(U;\mathbb R^n)\)，则 \(\int_U\operatorname{div}X=0\)。Fubini 定理因而给出
	\[
		\int_Uu\operatorname{div}X\,\mathrm dx
		=\int_{\mathbb R}\left(\int_{E_t\cap U}\operatorname{div}X\,\mathrm dx\right)\mathrm dt.
	\]
	对 \(|X|\le1\) 取上确界，得到
	\begin{equation}\label{g3:eq:bv-coarea-dual-bound}
		V(u;U)\le\int_{\mathbb R}P(E_t;U)\,\mathrm dt.
	\end{equation}

	第二步考虑 \(v\in C^\infty(U)\cap W^{1,1}(U)\)。记
	\[
		h(s)=\mathcal H^{n-1}(U\cap v^{-1}(s)).
	\]
	\cref{g3:thm:lipschitz-coarea-formula}的局部形式给出
	\(\int_\mathbb Rh(s)\,\mathrm ds=\int_U|\nabla v|\,\mathrm dx<\infty\)。这里的局部形式可将 \(U\) 分成可数个互不交的 Borel 块，每块包含于一个相对紧的球；在各球上用光滑截断把 \(v\) 延为全空间 Lipschitz 函数，再将公式求和。
	取光滑单调函数 \(H\)，满足 \(H=0\) 于 \(( -\infty,0]\)，\(H=1\) 于 \([1,\infty)\)。当 \(\varepsilon\downarrow0\) 时，
	\(H((v-t)/\varepsilon)\to\chi_{\{v>t\}}\) 于 \(L^1(U)\)。由变差的下半连续性、光滑链式法则及上述余面积公式，
	\[
	\begin{aligned}
		P(\{v>t\};U)
		&\le\liminf_{\varepsilon\downarrow0}\int_U
		\frac1\varepsilon H'\!\left(\frac{v-t}{\varepsilon}\right)|\nabla v|\,\mathrm dx\\
		&=\liminf_{\varepsilon\downarrow0}\int_t^{t+\varepsilon}
		\frac1\varepsilon H'\!\left(\frac{s-t}{\varepsilon}\right)h(s)\,\mathrm ds.
	\end{aligned}
	\]
	在 \(h\) 的每个 Lebesgue 点，最后的加权平均趋于 \(h(t)\)：权重积分为一，而误差不超过
	\(\|H'\|_\infty\varepsilon^{-1}\int_t^{t+\varepsilon}|h(s)-h(t)|\,\mathrm ds\)。因此
	\[
		\int_{\mathbb R}P(\{v>t\};U)\,\mathrm dt
		\le\int_U|\nabla v|\,\mathrm dx.
	\]

	第三步令 \(u\in BV(U)\)，由\cref{g3:thm:bv-strict-approximation}取 \(v_j\in C^\infty(U)\cap W^{1,1}(U)\)，使
	\[
		\|v_j-u\|_{L^1(U)}\longrightarrow0,\qquad
		\int_U|\nabla v_j|\,\mathrm dx\longrightarrow|Du|(U).
	\]
	逐点的区间长度恒等式给出
	\[
		\int_{\mathbb R}\|\chi_{\{v_j>t\}}-\chi_{\{u>t\}}\|_{L^1(U)}\,\mathrm dt
		=\|v_j-u\|_{L^1(U)}.
	\]
	选取右端可求和的子列，便知几乎所有 \(t\) 上水平集示性函数在 \(L^1(U)\) 中收敛。对这些 \(t\) 使用\cref{g3:thm:bv-lower-semicontinuity}，再用 Fatou 引理及第二步，得到
	\[
		\int_{\mathbb R}P(E_t;U)\,\mathrm dt
		\le\liminf_j\int_{\mathbb R}P(\{v_j>t\};U)\,\mathrm dt
		\le|Du|(U).
	\]
	结合第一步便得等式。若 \(V(u;U)=\infty\)，第一步已经迫使右端为无穷，因此扩展实数的等式也成立。

	用相对紧有界开集递增穷尽任意开集，变差的对偶定义与单调收敛定理将等式推广到所有开集。对局部 BV 函数，取一个可数穷尽即可同时保证几乎所有 \(E_t\) 在整个 \(\Omega\) 内局部有限周长。最后，与切片证明一样，对可数开集基先用对偶定义得到 \(t\mapsto|D\chi_{E_t}|(U)\) 的可测性，再用单调类得到一个可测测度族。其积分定义的局部有限测度在全部开集上等于 \(|Du|\)，因而在全部 Borel 集上相等。这证明第一个测度等式；第二个由有限周长集合的结构定理逐个水平代入得到。
\end{proof}

例如取 \(u=a\chi_E\)，\(a>0\)。只有 \(0<t<a\) 的水平集贡献周长，所以余面积公式给出 \(|Du|=a|D\chi_E|\)。即使通常的梯度在几乎处处为零，这部分变差仍完整地出现在水平集边界上。

\subsection{Euclid 等周不等式}
\label{g3:subsec:2105:03}
我们证明带最优常数的 Euclid 等周不等式。证明先处理光滑区域，再用水平集逼近推广到有限周长集合。光滑部分采用 Neumann 方程与下接触集的方法，其一般各向异性形式可见\cite[第4节，第2986--2988页]{CabreRosOtonSerra2016ABP}。

\begin{lemma}\label{g3:lem:smooth-set-perimeter-approximation}
	若 \(E\subset\mathbb R^n\) 满足 \(|E|<\infty\)、\(P(E)<\infty\)，则存在有界开集 \(F_j\)，其边界为光滑紧超曲面，并且
	\[
		|F_j\mathbin{\triangle}E|\longrightarrow0,
		\qquad P(F_j)\longrightarrow P(E).
	\]
	允许 \(F_j=\varnothing\)。当 \(n=1\) 时，光滑紧边界就是有限个端点。
\end{lemma}
\begin{proof}
	因为 \(\chi_E\in BV(\mathbb R^n)\)，\cref{g3:cor:bv-compact-smooth-approximation}给出 \(v_j\in C_c^\infty(\mathbb R^n)\)，使
	\[
		0\le v_j\le1,\qquad \delta_j:=\|v_j-\chi_E\|_1\to0,
		\qquad \int|\nabla v_j|\to P(E).
	\]
	取 \(\eta_j\downarrow0\)，使 \(\delta_j/\eta_j\to0\)，并舍去有限项以保证 \(\eta_j<1/2\)。由 BV 余面积公式，
	\[
		\int_{\eta_j}^{1-\eta_j}P(\{v_j>t\})\,\mathrm dt\le\int|\nabla v_j|.
	\]
	第二卷第14章第2节的光滑 Sard 定理说明，\(v_j\) 的临界值组成一维零测集。因此可以选取正则值 \(t_j\in(\eta_j,1-\eta_j)\)，满足
	\[
		P(\{v_j>t_j\})\le\frac1{1-2\eta_j}\int|\nabla v_j|+\frac1j.
	\]
	令 \(F_j=\{v_j>t_j\}\)。正则值定理给出光滑边界；因 \(t_j>0\) 且 \(v_j\) 紧支，\(F_j\) 及其边界有界。对 \(F_j\mathbin{\triangle}E\) 中的几乎每一点，\(|v_j-\chi_E|\ge\eta_j\)，故
	\[
		|F_j\mathbin{\triangle}E|\le\delta_j/\eta_j\longrightarrow0.
	\]
	上面的周长上界给出 \(\limsup P(F_j)\le P(E)\)，而下半连续性给出 \(P(E)\le\liminf P(F_j)\)，两者合起来完成证明。
\end{proof}

\begin{theorem}[Euclid 等周不等式]\label{g3:thm:sharp-euclidean-isoperimetric}
	设 \(n\ge2\)，\(E\subset\mathbb R^n\) 可测，\(|E|<\infty\)。则
	\begin{equation}\label{g3:eq:sharp-euclidean-isoperimetric}
		P(E)\ge n\omega_n^{1/n}|E|^{(n-1)/n}.
	\end{equation}
	球达到等号，故常数 \(n\omega_n^{1/n}\) 不能增大。当 \(n=1\)、\(0<|E|<\infty\) 时，相应结论为 \(P(E)\ge2\)，区间达到等号。
\end{theorem}
\begin{proof}
	先设 \(n\ge2\)，且 \(E=D\) 为有界、连通、非空的光滑区域。记
	\[
		V=|D|,\qquad S=P(D)=\mathcal H^{n-1}(\partial D),\qquad c=S/V.
	\]
	考虑 Neumann 问题
	\begin{equation}\label{g3:eq:isoperimetric-neumann}
		\operatorname{div}\nabla v=c\quad\text{于 }D,
		\qquad \partial_\nu v=1\quad\text{于 }\partial D.
	\end{equation}
	本书的非负 Laplace 算子为 \(\Delta=-\operatorname{div}\nabla\)，所以这里的内部方程写成 \(\Delta v=-c\)。\cref{g3:thm:weak-general-neumann}中的相容条件恰为
	\[
		\int_D(-c)\,\mathrm dx+\int_{\partial D}1\,\mathrm d\mathcal H^{n-1}=-cV+S=0.
	\]
	因此存在弱解，规定均值为零后唯一；\cref{g3:thm:boundary-higher-regularity}逐阶提升正则性，给出 \(v\in C^\infty(\overline D)\)。以下只使用它的二阶连续可微性。

	定义\term[xia jiechu ji]{下接触集}[lower contact set]
	\[
		\Gamma=\{x\in D:v(z)\ge v(x)+\nabla v(x)\cdot(z-x)
		\text{ 对所有 }z\in\overline D\}.
	\]
	它由那些切平面从下方支撑整张函数图像的点组成。对每个固定 \(z\)，定义中的不等式关于 \(x\) 为闭条件，故 \(\Gamma\) 在 \(D\) 中相对闭，特别地是 Borel 集。

	我们证明
	\begin{equation}\label{g3:eq:isoperimetric-gradient-cover}
		B_1(0)\subset\nabla v(\Gamma).
	\end{equation}
	固定 \(|p|<1\)，函数 \(z\mapsto v(z)-p\cdot z\) 在紧集 \(\overline D\) 上取得最小值。最小点不能在边界：在那里其外法向导数为
	\[
		\partial_\nu(v-p\cdot z)=1-p\cdot\nu\ge1-|p|>0,
	\]
	沿内法向稍作移动会严格降低函数值。于是最小点 \(x\) 位于 \(D\) 内，满足 \(\nabla v(x)=p\)，而全局极小条件正是 \(x\in\Gamma\)。覆盖关系由此成立。

	对 \(x\in\Gamma\)，函数 \(z\mapsto v(z)-v(x)-\nabla v(x)\cdot(z-x)\) 在 \(x\) 有局部极小值，所以 Hessian 矩阵 \(\nabla^2v(x)\) 半正定。若其特征值为 \(\lambda_1,\ldots,\lambda_n\ge0\)，算术--几何平均不等式给出
	\[
		0\le\det\nabla^2v(x)=\prod_{i=1}^n\lambda_i
		\le\left(\frac{\sum_i\lambda_i}{n}\right)^n
		=\left(\frac cn\right)^n.
	\]
	对梯度映射使用\cref{g3:thm:lipschitz-area-formula}，得到
	\[
	\begin{aligned}
		\omega_n
		&\le\mathcal L^n(\nabla v(\Gamma))
		\le\int_\Gamma|\det\nabla^2v|\,\mathrm dx\\
		&\le\left(\frac{S}{nV}\right)^n|\Gamma|
		\le\left(\frac{S}{nV}\right)^nV.
	\end{aligned}
	\]
	面积公式中的第二个不等式允许梯度映射有多个原像。它也适用于这里的局部 Lipschitz 映射：把 \(D\) 分成可数个互不交的 Borel 块，每块包含于相对紧的球，在各块上应用公式再求和，所得带重数积分至少覆盖像一次。整理最后一式即得
	\(S\ge n\omega_n^{1/n}V^{(n-1)/n}\)。

	若光滑有界区域 \(F\) 不连通，对其各连通分支 \(D_i\) 分别使用这个结论。光滑紧边界有有限个连通分支，而每个 \(D_i\) 都有非空边界，所以这里只有有限项；边界局部的半空间表示也保证一个边界分支只邻接一个内部连通分支。写 \(q=(n-1)/n\)，利用 \(a^q+b^q\ge(a+b)^q\)，得到
	\[
		P(F)=\sum_iP(D_i)
		\ge n\omega_n^{1/n}\sum_i|D_i|^q
		\ge n\omega_n^{1/n}|F|^q.
	\]
	其中二项幂次不等式可由 \((1+t)^q\le1+t^q\) 的导数比较得到，再作归纳。

	一般 \(E\) 若周长无穷，结论直接成立；若体积为零，则 \(P(E)=0\)。其余情形使用\cref{g3:lem:smooth-set-perimeter-approximation}，在光滑集合的不等式中令 \(j\to\infty\)，便得到式\eqref{g3:eq:sharp-euclidean-isoperimetric}。对半径 \(r\) 的球，\(|B_r|=\omega_nr^n\)、\(P(B_r)=n\omega_nr^{n-1}\)，代入即为等号。

	最后设 \(n=1\)。若 \(P(E)<\infty\)，\cref{g3:prop:one-dimensional-perimeter-structure}说明 \(E\) 模零测集是有限个互不相交区间的并。有限正长度排除了无界区间和空集，每个非空区间贡献两个端点，故 \(P(E)\ge2\)。一个有界非退化区间恰有两个端点，所以等号成立。
\end{proof}

这个证明中，边界条件使所有长度小于一的斜率都出现于下接触集上；内部方程控制 Hessian 的迹；面积公式再把梯度像的体积转换成 Hessian 行列式的积分。三者共同产生精确的常数 \(n\omega_n^{1/n}\)。

\subsection{固定体积的周长极小化}
\label{g3:subsec:2105:04}
对给定 \(v>0\)，定义全空间的等周函数
\[
	I_n(v)=\inf\{P(E):E\subset\mathbb R^n\text{ 可测},\ |E|=v\}.
\]
竞争者可以有不光滑的边界、多个连通分支，甚至仅以几乎处处等价类给出。周长的测度定义使这些集合都能进入同一个极小化问题。

\begin{proposition}\label{g3:prop:fixed-volume-ball-minimizer}
	设 \(n\ge1\)、\(v>0\)，定义 \(I_n(v)=\inf\{P(E):E\subset\mathbb R^n\text{ 可测},\ |E|=v\}\)，其中 \(P\) 为全空间周长、\(|E|\) 为 Lebesgue 体积、\(\omega_n\) 为单位球体积。若 \(n\ge2\)，则
	\[
		I_n(v)=n\omega_n^{1/n}v^{(n-1)/n}.
	\]
	任意中心、半径 \(r=(v/\omega_n)^{1/n}\) 的球都是极小集合。若 \(n=1\)，则 \(I_1(v)=2\)，而全部极小集合模零测集恰为长度 \(v\) 的区间。
\end{proposition}
\begin{proof}
	当 \(n\ge2\) 时，等周不等式对每个竞争者给出所示下界。半径为 \(r\) 的球具有体积 \(v\) 与周长 \(n\omega_nr^{n-1}\)，达到这个下界。平移保持体积与周长，所以球心任意。

	当 \(n=1\) 时，上一证明已经给出下界二，长度 \(v\) 的区间达到下界。反过来，极小集合的周长为二，一维结构定理将它表示为有限个有界区间的并，每个分支贡献两个端点，故恰有一个区间；体积条件规定其长度为 \(v\)。
\end{proof}

尺度律解释了等周函数中的幂次。对 \(\lambda>0\)，变量代换及周长的对偶定义给出
\[
	|\lambda E|=\lambda^n|E|,\qquad P(\lambda E)=\lambda^{n-1}P(E).
\]
为验证第二式，将测试场写成 \(Y(x)=X(\lambda x)\)，则 \(|Y|\le1\) 与 \(|X|\le1\) 等价，而且
\[
	\int_{\lambda E}\operatorname{div}X\,\mathrm dz
	=\lambda^{n-1}\int_E\operatorname{div}Y\,\mathrm dx.
\]
两边对测试场取上确界即可。因此只要某个形状达到极小，将它缩放到另一个体积仍达到相应极小值。

\begin{example}\label{g3:exm:fragmented-balls}
	设整数 \(n\ge2\)、\(k\ge1\) 及实数 \(v>0\)。在 \(\mathbb R^n\) 中将总体积 \(v\) 分给 \(k\) 个闭包两两不交的等半径球。每个球的体积为 \(v/k\)，各边界的周长相加，得到
	\[
		P(E)=k\,n\omega_n^{1/n}(v/k)^{(n-1)/n}
		=k^{1/n}I_n(v).
	\]
	只要 \(k>1\)，这个周长就严格大于一个同总体积球的周长。对于体积分配不相等的两个非空球，严格的幂次不等式同样给出周长增加。
\end{example}

若将竞争者限制在一个固定紧容器中，\cref{g3:thm:finite-perimeter-direct-method}给出了非空约束类的极小集合存在性。全空间模型则有平移自由：一列球可以整体移向无穷远，保持体积与周长，却在每个固定有界区域内趋于空集。这里球的显式构造与尖锐下界直接给出全局极小值，因而不需要先从任意极小化序列中恢复一个整体极限。
% 纲要条目排为编号小节标题，并进入目录与书签；说明段仍保留于纲要注释中。
